\documentclass{article}
\usepackage[utf8]{inputenc}
\usepackage{amsmath}
\usepackage{amssymb}
\usepackage{graphicx}
\usepackage{xltxtra}
\usepackage{fontspec}
\setmainfont{Takao PMincho} % Or your preferred Japanese font

\title{A LL(1) Parser for a Fragment of Classical Japanese}
\author{J. Kawasaki}
\date{\today}

\begin{document}

\maketitle

\begin{abstract}
This paper presents the design and implementation of a parser for a defined fragment of Classical Japanese. We formulate an LL(1) grammar that covers basic sentence structures, including subordinate clauses, and implement a corresponding recursive-descent parser in the Go programming language. The parser constructs an Abstract Syntax Tree (AST), providing a structured representation of the input text suitable for further linguistic analysis or processing. This work serves as a practical, end-to-end example of applying compiler construction techniques to a natural language, offering a clear and reproducible case study.
\end{abstract}

\section{Introduction}
The formal analysis of natural language syntax has long benefited from the tools and methodologies of computer science, particularly from the theory of formal languages and compilers. This paper details a project to build a parser for a limited subset of Classical Japanese (Yamato-kotoba). The primary goal is pedagogical: to provide a concrete example of how a formal grammar can be defined for a natural language fragment and how a corresponding parser can be implemented from scratch.

This paper is structured as follows: Section 2 details the LL(1) grammar. Section 3 describes the implementation of the lexer and the recursive-descent parser in Go. Finally, Section 4 concludes and suggests directions for future work.

\section{Grammar}
The grammar was initially designed for a fragment of Classical Japanese and subsequently extended to handle basic modern structures.

\subsection{Initial Grammar (Classical Japanese)}
\begin{verbatim}
grammar ClassicalJapanese;

// --- Top Level ---
sentence: clause+ EOF;

// --- Clause and Phrase Structure ---
clause: np vp;

np: baseNP npTail;
npTail: relClause npTail | /* epsilon */;
baseNP: NOUN caseParticle?;

relClause: adjPhrase | clause;

// --- Verb Phrase and Auxiliaries ---
vp: (np)? VERB auxList caseParticle?;
auxList: AUX auxList | /* epsilon */;

// --- Adjective Phrase ---
adjPhrase: ADJ AUX?;

// --- Morphemes ---
caseParticle: 'が' | 'を' | 'に';
AUX: 'き' | 'む' | 'ず';

NOUN: [word];
VERB: [word];
ADJ: [word];
\end{verbatim}

\subsection{Extended Grammar (Modern Japanese)}
To parse modern Japanese sentences, we extended the grammar to include common particles ('は', 'の'), polite forms ('です', 'ます'), and basic tense ('た').

\begin{verbatim}
grammar ModernJapanese;

vp: (np)? (verbPhrase | copulaPhrase);
verbPhrase: VERB auxList;
copulaPhrase: (NOUN | ADJ) COPULA;

caseParticle: 'が' | 'を' | 'に' | 'は' | 'の';
AUX: 'ます' | 'た' | ... ;
COPULA: 'です';
\end{verbatim}

\subsection{Lexical Tokens}
% ...

\subsection{Syntactic Rules}
% ...

\section{Implementation}
The parser is implemented in Go, following a standard recursive-descent approach. The implementation is divided into three main components: the lexer, the parser, and the Abstract Syntax Tree (AST) definitions.

\subsection{Lexer}
The lexer (or tokenizer) processes the input string, a sequence of characters, and converts it into a sequence of tokens. It handles multi-byte UTF-8 characters, correctly identifying Japanese words as well as single-character particles and auxiliaries. Keywords such as 'が' or 'き' are identified as distinct token types, while other word-like sequences are initially classified as generic nouns.

\subsection{Parser and AST}
The parser is a top-down, recursive-descent parser. It takes the token stream from the lexer and attempts to build a valid parse tree according to the grammar. For each non-terminal symbol in the grammar (e.g., \texttt{clause}, \texttt{np}, \texttt{vp}), there is a corresponding function in the parser.

The result of a successful parse is an Abstract Syntax Tree (AST). The AST is a tree representation of the abstract syntactic structure of the source code. Each node of the tree denotes a construct occurring in the source code. The AST nodes are defined as Go structs that correspond to the grammar's rules, such as \texttt{Clause}, \texttt{NP}, and \texttt{VP}. This tree can then be used by subsequent phases, such as an interpreter or code generator.

The AST was extended with nodes like \texttt{CopulaPhrase} and \texttt{VerbPhrase} to represent these new structures. The parser was updated to distinguish between these phrases based on lookahead.

\section{Conclusion}
We have successfully designed and implemented an LL(1) parser for a small but non-trivial fragment of Classical Japanese. The project demonstrates the viability of using standard parsing techniques to analyze the structure of a natural language, producing a well-defined Abstract Syntax Tree.

We successfully extended the parser to cover basic modern Japanese sentences, including polite forms. This demonstrates the extensibility of the architecture. Further work can incorporate more complex structures like passive and causative verbs.

\end{document} 